\documentclass{article}
\usepackage{amsmath,amssymb,xurl}
\usepackage[hidelinks]{hyperref}
% Shared commands for notes in docs/paper-gaps/.

\DeclareUnicodeCharacter{2080}{\ifmmode{}_{0}\else\textsubscript{0}\fi}
\DeclareUnicodeCharacter{2081}{\ifmmode{}_{1}\else\textsubscript{1}\fi}
\DeclareUnicodeCharacter{2082}{\ifmmode{}_{2}\else\textsubscript{2}\fi}
\DeclareUnicodeCharacter{2099}{\ifmmode{}_{n}\else\textsubscript{n}\fi}

\newcommand{\Lean}{\textsc{Lean}}
\ExplSyntaxOn
\NewDocumentCommand{\leanid}{m}{
  \ifmmode
    \textsf{\detokenize{#1}}
  \else
    \group_begin:
    \urlstyle{sf}
    \tl_set:Nn \l_tmpa_tl {#1}
    \tl_replace_all:Nnn \l_tmpa_tl {\_} {_}
    \exp_args:NV \path \l_tmpa_tl
    \group_end:
  \fi
}
\ExplSyntaxOff
\providecommand{\tr}{\operatorname{tr}}
\newcommand{\tnleanIssueBase}{https://github.com/LionSR/TNLean/issues/}
\newcommand{\tnleanPrBase}{https://github.com/LionSR/TNLean/pull/}
\newcommand{\tnleanDocBase}{https://sirui-lu.com/TNLean/docs/}
\newcommand{\ghissue}[1]{\href{\tnleanIssueBase#1}{GitHub issue \##1}}
\newcommand{\ghpr}[1]{\href{\tnleanPrBase#1}{PR \##1}}
\newcommand{\doclink}[1]{\href{\tnleanDocBase#1.html}{\path{#1}}}

% Dirac notation and the identity operator, as in blueprint/src/macros/common.tex.
% \providecommand keeps note-local definitions authoritative.
\usepackage{dsfont}
\providecommand{\ket}[1]{|#1\rangle}
\providecommand{\bra}[1]{\langle #1|}
\providecommand{\braket}[2]{\langle #1 | #2 \rangle}
\providecommand{\ketbra}[2]{|#1\rangle\!\langle #2|}
\providecommand{\Id}{\mathds{1}}
\providecommand{\one}{\mathds{1}}
\providecommand{\C}{\mathbb{C}}
\providecommand{\MN}[1]{M_{#1}(\C)}
\providecommand{\MD}{M_D(\C)}

% Verdict marker: \gapnote{<kind>}{<status>}. Kind is the policy
% classification (clarification, local-correction, scope-restriction,
% unfaithful, false-source, open-gap); status is open, wip (elimination
% underway), resolved, or historical. Machine-read by the site index;
% prints a verdict line.
\newcommand{\gapnote}[2]{\par\noindent\textbf{Verdict:} #1 (#2).\par}

\title{Restricted Dual Homotopy Needs More Than Winding}
\author{}
\date{}
\begin{document}
\maketitle
\gapnote{scope-restriction}{open}

Schuch, Cirac, and P\'erez-Garc\'ia,
\href{https://arxiv.org/abs/1001.3807v3}{arXiv:1001.3807v3},
Lemma 6.14, source \path{Papers/1001.3807/paper_v3.tex}, lines 2199--2214,
uses virtual invariance to deform a connecting flux string. The generated
Lean homotopy records the primal vertices allowed to support square moves.
Completeness on the full torus does not imply completeness after some
square moves are forbidden.

The definition \leanid{TNLean.PEPS.torusDualSeamWord} reads a path's ordered
seam crossings in an arbitrary test group, assigning $u,v$ to its eastward
and northward seam crossings and inverses to reverse crossings.
The theorem
\leanid{TNLean.PEPS.TorusDualHomotopy.seamWord_eq_of_origin_not_mem}
proves invariance of this nonabelian word under every generated homotopy
whose swept support excludes the origin. Backtracks cancel, and a square
can interchange both seam generators only when its enclosed primal vertex
is the forbidden origin.

Consequently
\leanid{TNLean.PEPS.torusDualSquareLoop_obstructs_restricted_completeness}
proves that the square loop based at $(-1,-1)$ has zero lifted displacement
but cannot contract by moves away from the origin, using any noncommuting
pair $u,v$ in a test group. Its word is $uvu^{-1}v^{-1}\ne1$, whereas the
constant path has word $1$. Noncommuting transpositions in the symmetric
group on three elements give a concrete choice of test group. The Lean
theorem is stated for every such noncommuting pair and all torus periods.
The unconditional version
\leanid{TNLean.PEPS.torusDualSquareLoop_not_homotopic_avoiding_origin}
constructs that pair internally and proves noncommutation by finite
kernel-checked computation. No noncommuting witness or physical flux group
is supplied to this version.

Thus equality of lifted displacement is insufficient for restricted-support
homotopy completeness. A positive classification must retain the additional
nonabelian information arising from forbidden cells. Neither this invariant
nor the obstruction theorem asserts such completeness. The seam word is
not a physical flux insertion or a statement that the two physical PEPS
states differ. In particular, full-support virtual-invariance arguments
can still identify physical coefficient vectors. Restricted completeness
and general noncentral open-pair density remain under \ghissue{8703}.
The full-lattice Hamiltonian scope of \ghissue{8947} is unchanged.
\end{document}
